% !TEX root = ../ApproxDA-TransUNet.tex
\section{Introduction}
\label{sec:introduction}

\IEEEPARstart{V}{ision} Transformers (ViTs) have fundamentally transformed medical image segmentation by capturing global long-range contextual relationships that convolutional neural networks (CNNs) struggle to resolve due to their localized receptive fields~\cite{chen2021transunet,cao2022swinunet,hatamizadeh2022unetr}. Within hybrid CNN-Transformer architectures, \emph{dual attention}~\cite{fu2019dual}---which concurrently models spatial dependencies via a Position Attention Module (PAM) and cross-channel interactions via a Channel Attention Module (CAM)---has emerged as a particularly potent paradigm. Pioneering frameworks such as DA-TransUNet~\cite{sun2024datransunet} embedded dual-attention blocks across multi-scale skip connections, establishing benchmark performance on multi-organ and polyp segmentation.

However, the computational cost of full-map dual attention scales quadratically with spatial token length ($\mathcal{O}(N^2C)$) and channel depth ($\mathcal{O}(C^2N)$). To mitigate this bottleneck, a rich body of efficient approximation strategies has emerged, including local windowed partitioning~\cite{liu2021swin}, low-rank matrix decomposition~\cite{wang2020linformer}, and grouped channel factorizations~\cite{rahman2024emcad}. A ubiquitous, yet rarely scrutinized, assumption underlying these approximation techniques is that sparse or localized attention serves purely as an engineering compromise: a computational compression step intended to approximate global interactions with negligible loss in representational fidelity.

In medical vision, however, this assumption encounters a profound challenge: biomedical segmentation tasks exhibit extreme anatomical and contextual heterogeneity. While segmenting the human pancreas in abdominal CT requires resolving spatial relationships relative to distant anatomical landmarks (such as the stomach, duodenum, and splenic vein), segmenting a localized colon polyp or cutaneous melanoma depends predominantly on fine-grained textural cues and local boundary gradients. This dichotomy raises a fundamental research question: \emph{How does attention approximation alter representation learning across heterogeneous clinical tasks, and what intrinsic task properties govern when and at what scale approximation is beneficial?}

To address this question, we develop \textbf{ApproxDA-TransUNet}\footnote{Source code and pretrained weights will be made available upon acceptance.}, a principled dual-attention framework featuring three orthogonal, independently controllable approximation axes: (i) spatial window scope ($M$), (ii) representation projection rank ($r$), and (iii) channel grouping count ($G$). Unlike conventional lightweight networks optimized solely for FLOP reduction, ApproxDA-TransUNet is engineered with clean modular decoupling, enabling systematic isolation of individual approximation factors without architectural confounding. Furthermore, by restructuring tensor operations for contiguous memory layout, ApproxDA-TransUNet avoids a PyTorch autograd/DDP incompatibility that we encountered in our reproduced DA-TransUNet implementation, enabling native multi-GPU Distributed Data Parallel (DDP) training in our test environment.

Systematic evaluation across five clinical benchmarks spanning four imaging modalities (Synapse multi-organ CT, ACDC cardiac MRI, Kvasir-SEG and CVC-ClinicDB endoscopic/colonoscopic polyp imaging, and ISIC 2018 dermoscopy) reveals a structured spectrum of context sensitivity. Building on these observations, we conduct a systematic candidate-factor analysis to identify the property most consistently associated with attention window sensitivity, analyze late-stage optimization dynamics, and provide a gradient-based fixed-point analysis for multi-branch gating behavior.

The primary contributions of this paper are summarized as follows:
\begin{itemize}
  \item \textbf{Modular ApproxDA-TransUNet Architecture with DDP-Compatible Implementation:} We propose a modular dual-attention framework with three independently controllable approximation axes ($M, r, G$) that outperforms DA-TransUNet in Dice across all five benchmarks while remaining competitive with representative CNN- and Transformer-based methods, and which avoided a DDP incompatibility we encountered with our reproduced DA-TransUNet baseline in our test environment.
  \item \textbf{Discovery of GCS and the SSD Hypothesis:} Across five benchmarks spanning four imaging modalities, we observe a $4.6\times$ difference in single-run window-sensitivity spans ($\Delta\text{DSC}$ span: $0.50$\,pp to $2.30$\,pp), formalizing this phenomenon as \emph{Global Context Sensitivity (GCS)}. Using ACDC as a key discriminating benchmark, we evaluate five candidate spatial metrics and find that inter-class window separation (SC5, $\rho{=}+0.89$) is the candidate most consistently associated with GCS across all evaluated benchmarks---an interpretive finding we term \emph{Semantic Spatial Dependency (SSD)}.
  \item \textbf{Optimization Dynamics \& Gate Collapse Analysis:} We characterize task-dependent differences in late-stage validation variability across the selected ApproxDA configurations, while explicitly accounting for differences in gate mode and window size (up to $18.2\times$ lower variability on Kvasir-SEG and ISIC 2018, but a reverse trend on Synapse). We further provide a gradient-based fixed-point analysis showing that symmetric dual-branch architectures tend toward uniform gating ($g \approx 0.5$), supported by empirical gate measurements across training.
  \item \textbf{Exploratory Design Heuristics:} We discuss practical heuristics for attention architecture design based on our findings, including pre-training SC5 assessment as a lightweight indicator of expected window sensitivity.
\end{itemize}

The remainder of this paper is organized as follows: Section~\ref{sec:related_work} reviews related literature; Section~\ref{sec:method} presents the ApproxDA-TransUNet framework and complexity formulations; Section~\ref{sec:experiments} reports results against state-of-the-art baselines across five benchmarks; Section~\ref{sec:gcs_mechanism} investigates the GCS spectrum and evaluates candidate mechanisms including Semantic Spatial Dependency (SSD); Section~\ref{sec:generalization} analyzes training optimization volatility and gate collapse dynamics; Section~\ref{sec:conclusion} concludes with practical design guidelines.
